\par\Needspace{6\baselineskip}\medskip
\SourceLeaf{213}\hypertarget{native-213-article157}{}\pdfbookmark[2]{157.}{article-157}
\begin{center}157.\end{center}
Zunächst bestimmen wir nun die Lage der Großkreise, die von den heliozentrischen Orten der Erde zu den geozentrischen Orten des Planeten gezogen werden. Die Schnittpunkte dieser Kreise mit dem Äquator oder, wenn man will, ihre aufsteigenden Knoten denken wir uns mit den Zeichen $\mathfrak A,\mathfrak A',\mathfrak A''$ versehen. Die Abstände der Punkte $B,B',B''$ von diesen Schnittpunkten bezeichnen wir mit $\Delta,\Delta',\Delta''$. Bei den meisten Rechenschritten müssen wir nun $A,A',A''$ durch $\mathfrak A,\mathfrak A',\mathfrak A''$ und $\delta,\delta',\delta''$ durch $\Delta,\Delta',\Delta''$ ersetzen. Wo hingegen $A,A',A'',\delta,\delta',\delta''$ beizubehalten sind, wird der aufmerksame Leser auch ohne unseren Hinweis leicht erkennen.

Die Rechnung ergibt nun
\begin{center}
\begingroup\footnotesize\setlength{\tabcolsep}{3pt}\renewcommand{\arraystretch}{1.25}
\begin{tabular}{l|r|r|r}
\shortstack{Rektasz.\ der Punkte\\$\mathfrak A,\mathfrak A',\mathfrak A''$} & $233^0\,54'\,57''{,}10$ & $253^0\,8'\,57''{,}01$ & $276^0\,40'\,25''{,}87$\\
$\gamma,\gamma',\gamma''$ & $51\quad17\quad15{,}74$ & $90\quad1\quad3{,}19$ & $131\quad59\quad58{,}03$\\
$\Delta,\Delta',\Delta''$ & $215\quad58\quad49{,}27$ & $212\quad52\quad48{,}96$ & $220\quad9\quad12{,}96$\\
$\delta,\delta',\delta''$ & $56\quad26\quad34{,}19$ & $55\quad26\quad31{,}79$ & $69\quad10\quad57{,}84$\\
$\mathfrak A'D,\mathfrak AD',\mathfrak AD''$ & $23\quad54\quad52{,}13$ & $30\quad18\quad3{,}25$ & $29\quad8\quad43{,}32$\\
$\mathfrak A''D,\mathfrak A''D',\mathfrak A'D''$ & $33\quad3\quad26{,}35$ & $31\quad59\quad21{,}14$ & $22\quad20\quad6{,}91$\\
$\epsilon,\epsilon',\epsilon''$ & $47\quad1\quad54{,}69$ & $89\quad34\quad57{,}07$ & $42\quad33\quad41{,}17$\\
Logarithmen der Sinuswerte & $9{,}864\,3525$ & $9{,}999\,9885$ & $9{,}830\,1910$\\
$\log\sin\tfrac12\epsilon'$ & & $9{,}847\,8971$ &\\
$\log\cos\tfrac12\epsilon'$ & & $9{,}851\,0614$ &
\end{tabular}\endgroup
\end{center}

In der Rechnung nach Artikel 138 wird für $l'$ die Rektaszension des Punktes $\mathfrak A'$ verwendet. So ergibt sich
\begin{center}\begin{tabular}{ll}
$\log T\sin t$ .\,.\,.\,. & $8{,}486\,8236_n$\\
$\log T\cos t$ .\,.\,. & $9{,}284\,8162_n$.
\end{tabular}\end{center}
\noindent Somit gilt $t=189^0\,2'\,48''{,}83$, $\log T=9{,}290\,2527$ und \mbox{$t+\gamma'=279^0\,3'\,52''{,}02$},
\begin{center}\begin{tabular}{ll}
$\log S$ .\,.\,.\,.\,.\,.\,.\,.\,.\,.\,. & $9{,}011\,0566_n$\\
$\log T\sin(t+\gamma')$ .\,.\,. & $9{,}284\,7950_n$,
\end{tabular}\end{center}
\noindent woraus sich $\Delta'-\sigma=208^0\,1'\,55''{,}64$ und $\sigma=4^0\,50'\,53''{,}32$ ergeben.

\Needspace{13\baselineskip}
In den Formeln aus Artikel 140 für $a$, $b$ und $\dfrac b a$ sind $\sin\delta$, $\sin\delta'$, $\sin\delta''$ selbst beizubehalten, ebenso in den Formeln aus Artikel 142. Für diese Rechnungen haben wir
\SourceLeaf{214}
\begin{center}\begingroup\footnotesize\setlength{\tabcolsep}{4pt}\renewcommand{\arraystretch}{1.3}
\begin{tabular}{ll}
$\mathfrak A''D'-\Delta''=171^0\,50'\,8''{,}18;$ & \shortstack{$\log\sin$ .\,.\,. $9{,}152\,3306;$\\$\log\cos$ .\,. $9{,}995\,5759_n$}\\
$\mathfrak AD'-\Delta=174\quad19\quad13{,}98;$ & \shortstack{--- .\,.\,. $8{,}995\,4722;$\\--- .\,.\,. $9{,}997\,8629_n$}\\
$\mathfrak A''D-\Delta''=172\quad54\quad13{,}39;$ & --- .\,.\,. $9{,}091\,7972$\\
$\mathfrak A'D-\Delta'+\sigma=175\quad52\quad56{,}49;$ & --- .\,.\,. $8{,}856\,1520$\\
$\mathfrak AD''-\Delta=173\quad9\quad54{,}05;$ & --- .\,.\,. $9{,}075\,5844$\\
$\mathfrak A'D''-\Delta'+\sigma=174\quad18\quad11{,}27;$ & --- .\,.\,. $8{,}996\,7978$.
\end{tabular}\endgroup\end{center}
\noindent Daraus erhalten wir
\[
\begin{aligned}
\log\varkappa&=0{,}921\,1850&\log\lambda&=0{,}081\,2057_n\\
\log\varkappa''&=0{,}811\,2762&\log\lambda''&=0{,}031\,9691_n\\
\log a&=0{,}109\,9088,&a&=+1{,}287\,9790\\
\log b&=0{,}181\,0404&&\\
\log\frac b a&=0{,}071\,1314&&
\end{aligned}
\]
\noindent woraus sich $\log b=0{,}181\,0402$ ergibt. Aus diesen beiden nahezu gleichen Werten übernehmen wir den Mittelwert $\log b=0{,}181\,0403$. Schließlich ergibt sich
\[
\begin{aligned}
\log c&=1{,}045\,0295\\
d&=+0{,}448\,9906\\
\log e&=9{,}210\,2894,
\end{aligned}
\]
\noindent Damit sind die Vorrechnungen abgeschlossen.

Das Zeitintervall zwischen der zweiten und dritten Beobachtung beträgt $39{,}874\,409$ Tage, zwischen der ersten und zweiten $30{,}900\,961$ Tage. Daraus ergibt sich $\log\theta=9{,}836\,2757$, $\log\theta''=9{,}725\,5533$. Wir setzen also \textit{für die erste Hypothese}
\[
\begin{aligned}
x=\log P&=9{,}889\,2776\\
y=\log Q&=9{,}561\,8290.
\end{aligned}
\]
\noindent Als wichtigste Ergebnisse der Rechnung erhält man dann:
\[
\begin{aligned}
\omega+\sigma&=20^0\,8'\,46''{,}72\\
\log Qc\sin\omega&=0{,}028\,2028.
\end{aligned}
\]
\noindent Daraus ergibt sich der wahre Wert $z=21^0\,11'\,24''{,}30$ sowie $\log r'=0{,}350\,9379$. Die drei übrigen Werte von $z$, die der Gleichung IV in Artikel 141 genügen, sind in diesem Fall
\[
\begin{aligned}
z&=63^0\,40'\,50''\\
z&=101\quad12\quad58\\
z&=199\quad24\quad7,
\end{aligned}
\]
\noindent\SourceLeaf{215}von denen der erste als Näherung an die Erdbahn anzusehen ist. Wegen des zu großen Zeitintervalls ist seine Abweichung hier allerdings viel größer als im vorhergehenden Beispiel. --- Die weitere Rechnung liefert die folgenden Zahlen:
\begin{center}\begingroup\footnotesize\setlength{\tabcolsep}{6pt}\renewcommand{\arraystretch}{1.3}
\begin{tabular}{ll}
$\zeta=195^0\,12'\,2''{,}48$ & $\zeta''=196^0\,57'\,50''{,}78$\\
$\log r=0{,}364\,7022$ & $\log r''=0{,}335\,5758$\\
$\tfrac12(u''+u)=266^0\,47'\,50''{,}47$ & $\tfrac12(u''-u)=-43^0\,39'\,5''{,}33$
\end{tabular}\endgroup\end{center}
\[
\begin{aligned}
2f'&=22^0\,32'\,40''{,}86\\
2f&=13\quad5\quad41{,}17\\
2f''&=9\quad27\quad0{,}05.
\end{aligned}
\]
\noindent Die Differenz zwischen $2f'$ und $2f+2f''$, die hier $0''{,}36$ beträgt, verteilen wir so auf $2f$ und $2f''$, dass wir $2f=13^0\,5'\,40''{,}96$, $2f''=9^0\,26'\,59''{,}90$ setzen.

Nun sind die Zeiten wegen der Aberration zu korrigieren. Dabei ist in den Formeln aus Artikel 145 $AD'-\zeta=\mathfrak AD'-\Delta+\delta-\zeta$, $A''D'-\zeta''=\mathfrak A''D'-\Delta''+\delta''-\zeta''$ zu setzen. Wir haben also
\begin{center}\begin{tabular}{ll}
$\log r$ .\,.\,.\,.\,.\,.\,.\,.\,.\,.\,.\,. & $0{,}36470$\\
$\log\sin(AD'-\zeta)$ .\,.\,. & $9{,}76462$\\
$\mathrm C.\log\sin\delta$ .\,.\,.\,.\,.\,.\,. & $0{,}07918$\\
$\log\mathrm{Konst.}$ .\,.\,.\,.\,.\,.\,.\,. & $7{,}75633$\\\hline
& $7{,}96483$\\
Zeitreduktion$=$ & $0{,}009222$
\end{tabular}\end{center}
\begin{center}\begin{tabular}{ll}
$\log r'$ .\,.\,.\,.\,.\,. & $0{,}35094$\\
$\log\sin(\delta'-z)$ & $9{,}75038$\\
$\mathrm C.\log\sin\delta'$ . & $0{,}08431$\\
& $7{,}75633$\\\hline
& $7{,}94196$\\
& $0{,}008749$
\end{tabular}\end{center}
\begin{center}\begin{tabular}{ll}
$\log r''$ .\,.\,.\,.\,.\,.\,.\,.\,.\,. & $0{,}33558$\\
$\log\sin(A''D'-\zeta'')$ & $9{,}84220$\\
$\mathrm C.\log\sin\delta''$ .\,.\,.\,.\,. & $0{,}02932$\\
& $7{,}75633$\\\hline
& $7{,}96343$\\
& $0{,}009192$
\end{tabular}\end{center}
\Needspace{8\baselineskip}
\noindent Daraus ergibt sich
\begin{center}\begin{tabular}{r|r|r}
Korrigierte Zeiten & Intervalle & Logarithmen\\
Nov.\ $5{,}564\,852$ & &\\
& $30{,}901\,434$ & $1{,}489\,9785$\\
$36{,}466\,286$ & &\\
& $39{,}873\,966$ & $1{,}60\overset{.}{0}\,6894$\\
$76{,}340\,252$ & &
\end{tabular}\end{center}
\noindent woraus sich die korrigierten Logarithmen der Größen $\theta,\theta''$ ableiten, und zwar $9{,}836\,2708$ beziehungsweise $9{,}725\,5599$. Beginnt man dann die Berechnung der Elemente aus $r',r'',2f,\theta$, so ergibt sich $\log\eta=0{,}003\,1921$, ebenso wie man aus $r,r',2f'',\theta''$ $\log\eta''=0{,}001\,7300$ erhält. Damit erhält man $\log P'=9{,}890\,7512$, $\log Q'=9{,}571\,2864$, also
\[X=+0{,}001\,4736,\qquad Y=+0{,}009\,4574.\]

Die wichtigsten Ergebnisse der \textit{zweiten Hypothese}, in der wir ansetzen
\SourceLeaf{216}
\[
\begin{aligned}
x=\log P&=9{,}890\,7512\\
y=\log Q&=9{,}571\,2864,
\end{aligned}
\]
\noindent sind folgende:
\begin{center}\begingroup\footnotesize\setlength{\tabcolsep}{6pt}\renewcommand{\arraystretch}{1.3}
\begin{tabular}{r|r}
$\omega+\sigma=20^0\,8'\,0''{,}87$ & $\log r''=0{,}336\,9708$\\
$\log Qc\sin\omega=0{,}037\,3071$ & $\tfrac12(u''+u)=267^0\,6'\,10''{,}75$\\
$z=21^0\,12'\,6''{,}09$ & $\tfrac12(u''-u)=-43\quad39\quad4{,}00$\\
$\log r'=0{,}350\,7110$ & $2f'=22\quad32\quad8{,}69$\\
$\zeta=195^0\,16'\,59''{,}90$ & $2f=13\quad1\quad54{,}65$\\
$\zeta''=196\quad52\quad40{,}63$ & $2f''=9\quad30\quad14{,}38$\\
$\log r=0{,}363\,0642$ &
\end{tabular}\endgroup\end{center}
\noindent Die Differenz $0''{,}34$ zwischen $2f'$ und $2f+2f''$ ist so zu verteilen, dass $2f=13^0\,1'\,54''{,}45$, $2f''=9^0\,30'\,14''{,}24$ angesetzt werden.

Wenn es sich zu lohnen scheint, die Zeitkorrekturen hier neu zu berechnen, so erhält man für die erste Beobachtung $0{,}009\,169$, für die zweite $0{,}008\,742$, für die dritte $0{,}009\,236$ und damit die korrigierten Zeiten Nov.\ $5{,}564\,905$, Nov.\ $36{,}466\,293$, Nov.\ $76{,}340\,208$. Daraus folgt
\[
\begin{aligned}
\log\theta&=9{,}836\,2703,&\log\theta''&=9{,}725\,5594\\
\log\eta&=0{,}003\,1790,&\log\eta''&=0{,}001\,7413\\
\log P'&=9{,}890\,7268,&\log Q'&=9{,}571\,0593.
\end{aligned}
\]
\noindent Auf diese Weise ergibt sich also aus der zweiten Hypothese
\[X=-0{,}000\,0244,\qquad Y=-0{,}000\,2271.\]

Schließlich setzen wir in der \textit{dritten Hypothese}
\[
\begin{aligned}
x=\log P&=9{,}890\,7268\\
y=\log Q&=9{,}571\,0593,
\end{aligned}
\]
\noindent Die wichtigsten Ergebnisse der Rechnung lauten:
\begin{center}\begingroup\footnotesize\setlength{\tabcolsep}{6pt}\renewcommand{\arraystretch}{1.3}
\begin{tabular}{r|r}
$\omega+\sigma=20^0\,8'\,1''{,}62$ & $\log r''=0{,}336\,9536$\\
$\log Qc\sin\omega=0{,}037\,0857$ & $\tfrac12(u''+u)=267^0\,5'\,53''{,}09$\\
$z=21^0\,12'\,4''{,}60$ & $\tfrac12(u''-u)=-43\quad39\quad4{,}19$\\
$\log r'=0{,}350\,7191$ & $2f'=22\quad32\quad7{,}67$\\
$\zeta=195^0\,16'\,54''{,}08$ & $2f=13\quad1\quad57{,}42$\\
$\zeta''=196\quad52\quad44{,}45$ & $2f''=9\quad30\quad10{,}55$\\
$\log r=0{,}363\,0960$ &
\end{tabular}\endgroup\end{center}
\noindent\SourceLeaf{217}Die Differenz $0''{,}30$ wird hier so verteilt, dass $2f=13^0\,1'\,57''{,}25$, $2f''=9^0\,30'\,10''{,}42$ angesetzt werden\footnote{Diese recht große und in allen Hypothesen nahezu gleiche Differenz entstand zum größten Teil dadurch, dass $\sigma$ um \underline{f}ast zwei Hundertstel einer Bogensekunde zu klein und der Logarithmus von $b$ um einige Einheiten zu groß ermittelt worden war.}.

Da sich all diese Zahlen nur sehr wenig von denen unterscheiden, die die zweite Hypothese geliefert hat, können wir nun sicher schließen, dass die dritte Hypothese keiner weiteren Korrektur bedarf und eine neue Hypothese somit überflüssig ist. Wir können daher jetzt zur Berechnung der Elemente aus $2f',\theta',r,r''$ übergehen. Da diese Rechnung aus den oben bereits ausführlich erläuterten Rechenschritten besteht, genügt es, hier für diejenigen, die sie selbst ausführen wollen, die daraus folgenden Elemente anzugeben:
\begin{center}\begingroup\footnotesize\setlength{\tabcolsep}{3pt}\renewcommand{\arraystretch}{1.3}
\begin{tabular}{ll}
\shortstack[l]{Rektaszension des aufsteigenden\\Knotens am Äquator\quad .\,.\,.\,.\,.\,.\,.\,.\,.\,.\,.\,.\,.\,.\,.\,.} & $158^0\,40'\,38''{,}93$\\
Neigung der Bahn gegen den Äquator\quad .\,.\,.\,.\,.\,.\,.\,.\,.\,.\,.\,.\,.\,.\,.\,.\,.\,.\,.\,.\,.\,.\,.\,.\,.\,.\,. & $11\quad42\quad49{,}13$\\
\shortstack[l]{Abstand des Perihels von diesem\\aufsteigenden Knoten\quad .\,.\,.\,.\,.\,.\,.\,.\,.\,.\,.\,.\,.\,.\,.\,.\,.\,.\,.} & $323\quad14\quad56{,}92$\\
Mittlere Anomalie für die Epoche 1806\quad .\,.\,.\,.\,.\,.\,.\,.\,.\,.\,.\,.\,.\,.\,.\,.\,.\,.\,.\,.\,.\,.\,.\,.\,. & $335\quad4\quad13{,}05$\\
Mittlere (siderische) tägliche Bewegung\quad .\,.\,.\,.\,.\,.\,.\,.\,.\,.\,.\,.\,.\,.\,.\,.\,.\,.\,.\,.\,.\,.\,.\,.\,.\,.\,.\,.\,.\,.\,. & $770''{,}2663$\\
Winkel $\varphi$\quad .\,.\,.\,.\,.\,.\,.\,.\,.\,.\,.\,.\,.\,.\,.\,.\,.\,.\,.\,.\,.\,.\,.\,.\,.\,.\,.\,.\,.\,.\,.\,.\,.\,.\,.\,.\,.\,.\,.\,.\,.\,.\,.\,.\,.\,. & $14^0\,9'\,3''{,}91$\\
Logarithmus der großen Halbachse\quad .\,.\,.\,.\,.\,.\,.\,.\,.\,.\,.\,.\,.\,.\,.\,.\,.\,.\,.\,.\,.\,.\,.\,.\,.\,.\,.\,.\,.\,.\,. & $0{,}442\,2438$.
\end{tabular}\endgroup\end{center}
